\documentclass[12pt,letterpaper]{article}
\usepackage[spanish]{babel}
\usepackage[utf8]{inputenc}
\usepackage[T1]{fontenc}
\usepackage[margin=2.5cm]{geometry}
\usepackage{setspace}
\usepackage{enumitem}
\usepackage{parskip}
\decimalpoint

\onehalfspacing

\begin{document}

\begin{center}
    {\Large \textbf{CONTRATO DE ARRENDAMIENTO DE INMUEBLE PARA VIVIENDA URBANA}}
\end{center}

\vspace{0.5cm}

Entre los suscritos a saber: \textbf{FERNANDO ARANDA OCHOA}, mayor de edad, identificado(a) con cédula de ciudadanía No. \textbf{13.718.122} expedida en \textbf{Bucaramanga Santander}, domiciliado(a) y residente en \textbf{Piedecuesta Cl 22B \#1W 76 Portal del Valle}, quien para efectos del presente contrato se denominará \textbf{EL ARRENDADOR}; y \textbf{AIDA YAMILE ORTEGA}, mayor de edad, identificado(a) con cédula de ciudadanía No. \textbf{1.095.918.586} expedida en \textbf{Giron Santander}, quien actúa en su propio nombre y representación y quien para efectos del presente contrato se denominará \textbf{EL ARRENDATARIO}, hemos convenido celebrar el presente \textbf{CONTRATO DE ARRENDAMIENTO DE VIVIENDA URBANA}, el cual se regirá por la Ley 820 de 2003 y demás normas concordantes, así como por las siguientes cláusulas:

\begin{enumerate}[label=\textbf{CLÁUSULA \Roman*.}, wide, labelindent=0pt, itemsep=0.4cm]

\item \textbf{OBJETO.} EL ARRENDADOR entrega a título de arrendamiento a EL ARRENDATARIO, y este recibe a satisfacción, el goce del \textbf{primer piso} del inmueble ubicado en \textbf{Piedecuesta Cl 22B \#1W 76 Portal del Valle}, en la ciudad de \textbf{Piedecuesta Santander}, con el fin de que sea destinado según lo dispuesto en la cláusula siguiente.

\item \textbf{DESTINACIÓN.} EL ARRENDATARIO destinará el inmueble arrendado exclusivamente para \textbf{VIVIENDA}, destinación que no podrá ser cambiada. En caso de incumplimiento, EL ARRENDADOR podrá dar por terminado el contrato y exigir la entrega inmediata del inmueble.

\textit{Parágrafo.} EL ARRENDATARIO no podrá destinar el inmueble a fines ilícitos, quedando prohibido su uso para ocultar o depositar armas, explosivos, sustancias estupefacientes, alucinógenas, inflamables o cualquier elemento que ponga en peligro la seguridad del inmueble o de sus vecinos.

\item \textbf{PRECIO Y FORMA DE PAGO.} El canon mensual de arrendamiento es la suma de \textbf{OCHOCIENTOS MIL PESOS} (\$\textbf{800.000}) M/CTE, que EL ARRENDATARIO pagará de manera anticipada dentro de los primeros \textbf{21} días de cada mes, en \textbf{Piedecuesta Santander}, durante toda la vigencia del contrato.

\textit{Parágrafo.} Como garantía del cumplimiento de las obligaciones derivadas del presente contrato, EL ARRENDATARIO entrega en este acto a EL ARRENDADOR la suma de \textbf{CUATROCIENTOS MIL PESOS} (\$\textbf{400.000}) a título de depósito, la cual será restituida al finalizar el contrato una vez verificado que el inmueble se entrega en las condiciones pactadas y a paz y salvo por todo concepto.

\item \textbf{ESTADO DEL INMUEBLE.} EL ARRENDATARIO declara recibir el inmueble en buen estado: paredes en buen estado y pintadas, vidrios y láminas sin roturas, baldosas en buen estado, instalaciones eléctricas en perfecto funcionamiento, llaves completas y el inmueble limpio. EL ARRENDATARIO se obliga a restituirlo en las mismas condiciones.

\item \textbf{TÉRMINO DEL CONTRATO.} El presente contrato tendrá una duración de \textbf{SEIS (6) MESES}, contados a partir de \textbf{los veintiún (21) días del mes de julio del año dos mil veintiséis (2026)}. Vencido este término, el contrato se prorrogará automáticamente por períodos iguales, salvo que alguna de las partes manifieste su voluntad de darlo por terminado conforme a lo dispuesto en la ley y en este contrato.


\item \textbf{OBLIGACIONES DEL ARRENDATARIO.} Además de las obligaciones legales, EL ARRENDATARIO se obliga a:
\begin{enumerate}[label=\alph*)]
    \item Pagar oportunamente el canon de arrendamiento.
    \item Conservar el inmueble en el mismo estado en que lo recibió.
    \item Informar oportunamente a EL ARRENDADOR sobre daños que requieran reparaciones necesarias, entendiendo por tales aquellas indispensables para el uso ordinario del inmueble, conforme al artículo 27 de la Ley 820 de 2003.
    \item Abstenerse de realizar mejoras o reformas sin autorización previa y escrita de EL ARRENDADOR.
    \item Abstenerse de ceder el contrato o subarrendar el inmueble, total o parcialmente, sin autorización previa y escrita de EL ARRENDADOR.
\end{enumerate}

\item \textbf{OBLIGACIONES DEL ARRENDADOR.} Además de las obligaciones legales, EL ARRENDADOR se obliga a:
\begin{enumerate}[label=\alph*)]
    \item Asumir el pago oportuno de los servicios públicos del inmueble, correspondientes a acueducto (agua), energía, gas y aseo (recolección de residuos), durante toda la vigencia del contrato.
    \item Garantizar a EL ARRENDATARIO el uso y goce pacífico del inmueble arrendado.
    \item Realizar las reparaciones necesarias del inmueble, de conformidad con el artículo 27 de la Ley 820 de 2003.
\end{enumerate}



\item \textbf{CAUSALES DE TERMINACIÓN.} EL ARRENDADOR podrá dar por terminado el presente contrato, además de las causales legales, por:
\begin{enumerate}[label=\arabic*)]
    \item El no pago oportuno del canon de arrendamiento.
    \item La suspensión o pérdida de los servicios públicos por causas imputables a EL ARRENDATARIO.
    \item El subarriendo, cesión o destinación distinta a la pactada, sin autorización.
    \item La afectación reiterada de la tranquilidad de los vecinos o el uso del inmueble para actividades ilícitas.
    \item La realización de mejoras, ampliaciones o modificaciones sin autorización de EL ARRENDADOR.
\end{enumerate}

\item \textbf{RESTITUCIÓN DEL INMUEBLE.} A la terminación del contrato, EL ARRENDATARIO deberá restituir el inmueble en las mismas condiciones en que lo recibió, a paz y salvo por concepto de cánones de arrendamiento. De existir obligaciones pendientes, EL ARRENDADOR podrá disponer del depósito de garantía para cubrirlas.

\item \textbf{MÉRITO EJECUTIVO.} El presente contrato, junto con los documentos que conforme a la ley correspondan, presta mérito ejecutivo para exigir el pago de los cánones adeudados, los daños ocasionados al inmueble, así como cualquier otra suma a cargo de EL ARRENDATARIO derivada de este contrato.

\item \textbf{PREAVISO.} Las partes se obligan a comunicar por escrito su intención de no prorrogar el contrato o de dar por terminada la relación contractual, con una antelación mínima de \textbf{UN (1) MES} al vencimiento del término pactado.

\end{enumerate}

En constancia de aceptación, las partes firman el presente contrato en la ciudad de \textbf{Piedecuesta Santander}, a los \textbf{27 días del mes de julio de 2026}.

\vspace{0.8cm}

\noindent
\begin{minipage}[t]{0.48\textwidth}
\rule{6cm}{0.4pt} \\
Firma \hspace{1.5cm} ARRENDADOR \\[0.3cm]
\rule{6cm}{0.4pt} \\
Nombre: \textbf{FERNANDO ARANDA OCHOA} \\
C.C. \textbf{13.718.122}
\end{minipage}
\hfill
\begin{minipage}[t]{0.48\textwidth}
\rule{6cm}{0.4pt} \\
Firma \hspace{1.5cm} ARRENDATARIO \\[0.3cm]
\rule{6cm}{0.4pt} \\
Nombre: \textbf{AIDA YAMILE ORTEGA} \\
C.C. \textbf{1.095.918.586}
\end{minipage}

\end{document}

% ============================================================
% DATOS PENDIENTES POR REEMPLAZAR (reemplazar cada marcador [ ]
% por la información real antes de firmar el contrato):
%
% FERNANDO ARANDA OCHOA                  -> Nombre completo del arrendador
% 13.718.122                      -> Número de cédula del arrendador
% Bucaramanga Santander     -> Ciudad de expedición de la cédula del arrendador
% Piedecuesta Cl 22B \#1W 76 Portal del Valle               -> Dirección de residencia del arrendador

% AIDA YAMILE ORTEGA                -> Nombre completo del arrendatario
% 1.095.918.586                    -> Número de cédula del arrendatario
% Giron Santander   -> Ciudad de expedición de la cédula del arrendatario
% Piedecuesta Cl 22B \#1W 76 Portal del Valle                 -> Dirección completa del inmueble arrendado

% Piedecuesta Santander                    -> Ciudad donde se ubica el inmueble
% OCHOCIENTOS MIL PESOS                        -> Valor del canon en letras (ej. CUATROCIENTOS MIL PESOS)
% 800.000               -> Valor del canon en números (ej. 400.000)
% 21                           -> Día del mes en que se paga el canon (ej. 5)

% Piedecuesta Santander                        -> Ciudad donde se realiza el pago
% CUATROCIENTOS MIL PESOS                     -> Valor del depósito de garantía en letras
% 400.000            -> Valor del depósito de garantía en números
% SEIS (6) MESES                  -> Duración del contrato (ej. DOCE (12) MESES)

% a los veintiún (21) días del mes de julio del año dos mil veintiséis (2026)                       -> Fecha de inicio del contrato
% UN (1) MES                    -> Tiempo mínimo de preaviso (ej. UN (1) MES)
% Piedecuesta Santander                       -> Ciudad donde se firma el contrato
% 27 días del mes de julio de 2026                        -> Fecha de la firma (ej. 23 días del mes de julio de 2026)
% ============================================================
